\section{Results}
\label{sec:results}
\subsection{Distribution reduces bursts, with overhead}
Table~\ref{tab:conference-baselines} isolates work placement and compares
native alternatives. Distributing the matched scalar implementation reduces
the typical hop peak from \StudyBatchPeak{} to
\StudyCorePeak\,$\mu$s, a factor of \StudyPlacementRatio. Aggregate
instrumented cost rises from \StudyBatchCost{} to
\StudyCoreCost\,ns/sample, an increase of about 18\%. The same FFT
arithmetic is being performed in smaller pieces; the result is less
concentrated execution rather than less total work.

\begin{table}[htbp]
\centering\small
\begin{tabular}{@{}lrr@{}}
\toprule
Analysis path & Cost & Peak [sessions] \\
\midrule
Matched batch & 68.3 & 52.2 [51.2--53.1] \\
Matched distributed & 80.3 & 6.67 [6.25--7.08] \\
PFFFT batch & 41.3 & 13.7 [13.1--14.3] \\
PFFFT hybrid & 49.8 & 8.15 [7.4--8.9] \\
vDSP batch & 43.6 & 17.1 [16.8--17.4] \\
vDSP hybrid & 59.4 & 8.92 [8.85--8.98] \\
FFTW batch & 41.3 & 21.8 [20.7--22.9] \\
FFTW hybrid & 46.6 & 12.5 [11.6--13.4] \\
\bottomrule
\end{tabular}
\caption{Scalar analysis, $N=4096$, $H=1024$, $D=64$, Hann, no smoothing. Cost is instrumented ns/sample; peak is the typical hop peak in $\mu$s. Brackets give the two session summaries. Each row contains four processes, 16,384 blocks, and 1,024 complete hops. No block exceeded the 1,333\,$\mu$s budget.}
\label{tab:conference-baselines}
\end{table}


All three native hybrids also lower the peak relative to their own batch
paths. PFFFT, for example, moves from \StudyPffftBatchPeak{} to
\StudyPffftHybridPeak\,$\mu$s while cost rises from
\StudyPffftBatchCost{} to \StudyPffftHybridCost\,ns/sample. Native hybrids
can therefore recover much of the burst reduction with lower aggregate cost
than the custom distributed transform. The fully distributed path has the
lowest typical peak at this scalar configuration, but it has the highest
cost. There is no single winner on both measures.

Across the campaign's three baseline lengths, four implementation families,
and two sessions, the distributed or hybrid member has the lower session
hop-peak summary in all 24 comparisons. These are related observations on one
machine, not 24 independent replications. Moreover, every typical peak in
Table~\ref{tab:conference-baselines} is small relative to the
1,333\,$\mu$s block budget. This experiment establishes the scheduling
effect; complete-host measurements determine whether the saved headroom is
substantial in a patch.

\subsection{A shorter schedule delivers earlier, larger bursts}
A native hybrid can finish within $H$, $H/2$, or $H/4$ sample calls, then
only capture input until the next frame endpoint. Table~\ref{tab:conference-horizons}
shows this tradeoff for four-channel vDSP analysis at a larger transform and
shorter hop. Reducing the completion interval from 240 to 60 calls cuts
endpoint-to-publication age from 4.98 to 1.23\,ms. The typical peak rises
from \StudyHorizonFullPeak{} to
\StudyHorizonQuarterPeak\,$\mu$s, approaching the batch peak of
125\,$\mu$s. The unchanged native transform remains indivisible; the
surrounding work becomes more densely packed.

\begin{table}[htbp]
\centering\small
\begin{tabular}{@{}lrrr@{}}
\toprule
Finish after & Age (ms) & Cost & Peak \\
\midrule
Batch & 0 & 544 & 125 \\
$H$ & 4.98 & 589 & 68.6 \\
$H/2$ & 2.48 & 594 & 95.4 \\
$H/4$ & 1.23 & 568 & 121 \\
\bottomrule
\end{tabular}
\caption{Four-channel vDSP analysis, $N=16384$, $H=240$, $D=64$, Hann. Age is endpoint-to-publication delay; cost is ns/sample for all four channels; peak is $\mu$s. Each row contains four processes and 3,840 blocks. Only the completion interval changes between hybrid rows.}
\label{tab:conference-horizons}
\end{table}


This configuration also shows a limit of the custom FFT. Its four-channel SIMD
path costs \StudySimdExtremeCost\,ns/sample with a
\StudySimdExtremePeak\,$\mu$s typical peak. The full-hop vDSP hybrid
costs \StudyHorizonFullCost\,ns/sample with a
\StudyHorizonFullPeak\,$\mu$s peak. Here the native hybrid is better on
both recorded measures. The benefit of interrupting every butterfly therefore
depends on transform size, hop, channel organization, and native kernel cost.

\subsection{The effect survives one host configuration}
With four aligned Fourier modules on one Rack engine thread, the distributed
implementation has a \StudyHostCorePeak\,$\mu$s typical hop peak,
compared with \StudyHostBatchPeak\,$\mu$s for native vDSP batch and
\StudyHostHybridPeak\,$\mu$s for the vDSP hybrid
(Table~\ref{tab:conference-host}). These occupy approximately 15\%, 67\%,
and 32\% of the block budget. Unlike the scalar comparison, the difference
is a substantial fraction of the time available to the rest of the patch.
The cost is higher aggregate processing: production uses
\StudyHostCoreCost\,ns/sample, versus \StudyHostBatchCost{} for batch and
\StudyHostHybridCost{} for hybrid. Both distributed paths also defer
publication by nearly the 30\,ms hop.

\begin{table}[htbp]
\centering\small
\begin{tabular}{@{}lrrr@{}}
\toprule
Analysis & Cost & Peak [sessions] & Misses \\
\midrule
Distributed & 2300 & 201 [197--206] & 18 \\
vDSP batch & 1730 & 890 [885--894] & 20 \\
vDSP hybrid & 2040 & 423 [417--428] & 38 \\
\bottomrule
\end{tabular}
\caption{Actual Rack engine: one thread, four aligned four-channel Fourier modules, $N=2048$, $H=1440$, $D=64$, flattop. Cost is block wall ns/sample for the complete engine; peak is $\mu$s. Brackets give session summaries. Misses are late synthetic releases across 16,384 blocks per row, including wake lateness, not audio underruns.}
\label{tab:conference-host}
\end{table}


With four engine threads and four aligned modules, the typical peaks converge
to about 100--115\,$\mu$s. Rack synchronizes its workers between
samples~\cite{rackengine}, but the experiment cannot isolate that cost from
core placement or power-state effects. The empty four-thread engine itself
consumes about three CPU-seconds per second of nominal audio, indicating that
its workers remain active between paced releases. The one-thread engine can
sleep. The contrast therefore changes more than available parallelism.

Alignment also matters. With sixteen modules on four threads, production's
typical peak is \StudyHostManyCorePeak\,$\mu$s, compared with
\StudyHostManyBatchPeak\,$\mu$s for aligned vDSP batch. Staggering the
batch modules lowers their peak to
\StudyHostStaggerBatchPeak\,$\mu$s. Phase distribution can reduce the
coincidence of complete transforms without making them resumable, although an
individual module cannot generally choose the phases of unrelated modules.

Lower typical peaks do not establish better deadline reliability. In the
one-thread, four-module comparison, production, vDSP batch, and hybrid miss
18, 20, and 38 synthetic releases respectively, across 16,384 blocks per row.
With sixteen aligned modules on four threads, production misses 19 releases
and vDSP batch misses none. Empty-engine controls also have late releases.
Wake delays, interruptions, and backlog remain relevant even when the normal
processing burst is much smaller than the deadline.

\subsection{Operating and numerical limits remain visible}
Direct paced calls to a default Fourier module have a typical peak of
\StudyModuleDefaultPeak\,$\mu$s. At $N=16384$ with a 5\,ms hop, the
four-port module reaches \StudyModuleExtremePeak\,$\mu$s; enabling both
smoothing modes raises this to \StudyModuleBothPeak\,$\mu$s and produces
196 compute-budget exceedances in 3,840 blocks. Scheduling distributes the
work but cannot make arbitrarily demanding settings affordable. No paced
native complete-module control was collected at this extreme configuration.

All numerical replays pass their declared acceptance policies. The largest
ordinary float relative vector errors are below $4.1\times10^{-5}$, against
a $3\times10^{-4}$ budget. These results include high overlap,
non-power-of-two hops, silence, and serialized lifecycle changes. Relative
error can reach one in the separately flagged long-decay tail; that region
is covered only by its absolute floor. The tests substantiate the implemented
analysis within those policies, not uniform relative accuracy for every weak
spectral component.
